\documentclass[10pt,a4paper]{amsart}
\usepackage[margin=19mm]{geometry}
\usepackage{amsmath,amssymb,mathtools}
\usepackage[scr=rsfs,cal=euler]{mathalfa}
\usepackage{xfrac}
\usepackage[unicode,hidelinks,bookmarks=false]{hyperref}
\newcommand{\ie}{i.e.\ }
\newcommand{\cf}{cf.\ }
\newcommand{\resp}{respectively\ }
\newcommand{\Subsection}{\subsection}
\providecommand{\nobreakdash}{\nobreak\hbox{-}}
\setlength{\emergencystretch}{2em}
\multlinegap0pt
\usepackage{bm}
\usepackage{bbm}
\usepackage{dsfont}
\usepackage{xspace}
\usepackage{cancel}
\usepackage{mathtools}
\usepackage{cleveref}
\usepackage{amsthm}
\makeatletter
\@ifundefined{proposition}{\newtheorem{proposition}{Proposition}}{}
\makeatother
\newcommand{\R}{\mathbb{R}}
\newcommand{\Z}{\mathbb{Z}}
\newcommand{\e}{\varepsilon}
\newcommand{\lln}{\gamma_\kappa}
\newcommand{\logq}{\eta_q}
\title[Lower tail large deviations of the stochastic six vertex mod]{Lower tail large deviations of the stochastic six vertex model}
\author{Sayan Das, Yuchen Liao, Matteo Mucciconi}
\date{Source: arXiv:2407.08530v2 (CC BY 4.0)}
\begin{document}
\maketitle

\noindent\emph{Source and licence.} Lower tail large deviations of the stochastic six vertex model. Version arXiv:2407.08530v2 (CC BY 4.0), distributed under the Creative Commons Attribution 4.0 International licence, as recorded in the arXiv licence metadata for this exact version and shown on the abstract page https://arxiv.org/abs/2407.08530v2. Full rights notice: licenses/arxiv-2407.08530-CC-BY-4.0.txt.\par\smallskip

\noindent\textit{Editorial scope.} Counted unit: the Proposition whose source label is \texttt{prop:equicont} in the article (source0.tex lines 1583--1588, source label \texttt{prop:equicont}), delivered together with the article's own complete original proof of it (source0.tex lines 1590--1629, one \texttt{proof} environment of 40 lines). No mathematics is written, repaired or re-derived: statement and proof are the article's own text line for line. Required context retained uncounted: p.energystab (lines 1289--1293); cor:tconvex (lines 1461--1466); prop:f\_limit\_g+T (lines 1503--1509). Every definition, display and label of those lines is retained unchanged. Omitted from this package and not redistributed: 1--1288, 1294--1460, 1467--1502, 1510--1582, 1589--1589, 1630--2517. Nothing inside a retained excerpt is omitted or altered: every retained body is delivered exactly as the article's lines are written, so no omission receipt of any kind accompanies this package. Citations: the retained excerpts carry no citation command at all, so no bibliography entry of the article is transcribed and no reference is redistributed. Reference adaptation: none was needed and none was made; every reference inside the delivered unit resolves inside it, so no unresolved reference or citation is produced. Printed slips kept literal: no printed word, symbol, hypothesis or number of the counted statement or of its proof was altered. Layout and class: the article's own class is \texttt{amsart} with packages including geometry, amsmath, amssymb, amsthm, amsfonts, hyperref, graphicx, color, colortbl, upgreek, euscript, hhline, mathbbol, mathtools; the wrapper is the lane's pinned \texttt{amsart} editorial wrapper at 10pt on A4, which restates the theorem-like environments the retained text needs and adds the article's own macro definitions that the retained text needs (\texttt{\textbackslash R}, \texttt{\textbackslash Z}, \texttt{\textbackslash e}, \texttt{\textbackslash lln}, \texttt{\textbackslash logq}), with extra wrapper packages (bm, bbm, dsfont, xspace, cancel, mathtools, cleveref). The delivered numbering is the unit's own. Assets: no retained span contains a figure environment, a graphics-inclusion command, a PDF-page inclusion, a table or a TikZ picture; no image file is copied into this package, the package declares no asset dependency and no image file is redistributed with it. Licence: version arXiv:2407.08530v2 of this article is distributed under the Creative Commons Attribution 4.0 International licence, as recorded in the arXiv licence metadata for this exact version and shown on the abstract page https://arxiv.org/abs/2407.08530v2. A full rights notice is delivered at licenses/arxiv-2407.08530-CC-BY-4.0.txt. Characters: every retained excerpt is pure ASCII, so the delivered engine, XeLaTeX with its default Latin Modern text font, reports no missing character.
\begin{proposition} \label{p.energystab} For large enough $s$ we have 
\begin{align}\label{f.para}
    \mathcal{F}_\kappa(s)=\logq \frac{(s-1)^2}{2}+\kappa \log \frac{1-aq}{1-a}.
\end{align}
\end{proposition}

\begin{proposition}
\label{cor:tconvex} For any $\e>0$, there exists $N_{\e} \in \Z_{\ge 1}$ such that for all $v_1,v_2\in \R$ and for all  $N\ge N_\e$ we have
\begin{align*}
\frac12\big(\Phi_{\kappa,N}(v_1)+\Phi_{\kappa,N}(v_2)\big)+\e \ge \Phi_{\kappa,N}\bigg(\frac{v_1+v_2}{2}-N^{-1/6}\bigg).
\end{align*}
\end{proposition}

\begin{proposition} \label{prop:f_limit_g+T}
				Let $x \in \mathbb{R}$ and define the function $g(x)=\frac{x^2}{2}\logq$. Then  we have
				\begin{equation}
                \label{eq:lim_inf_conv}
					\mathcal{F}_\kappa(s) = \lim_{N\to \infty}  \left( g \oplus \Phi_{\kappa,N} \right)(s):=\lim_{N\to \infty}\inf_{y \in \mathbb{R}} \left\{ g(y) + \Phi_{\kappa,N}(s-y) \right\}.
				\end{equation}
			\end{proposition}

\begin{proposition}\label{prop:equicont} $\Phi_{\kappa,N}$ is equicontinuous in the following sense. For any $\varepsilon>0$, there exists $\delta>0$ and $N_\varepsilon>0$ such that for all $N\ge N_\varepsilon$ we have
    $$
    |\Phi_{\kappa,N}(x)-\Phi_{\kappa,N}(y)|\le \varepsilon,
    $$ 
    for all $x,y\in[\lln,1]$ with $|x-y|\le \delta$.
\end{proposition}

\begin{proof} 
    Let $x>y$ and assume $x-y=\delta$. By the midpoint convexity stated in \Cref{cor:tconvex}, for any fixed $\varepsilon'$, there exists $N_{\varepsilon'}$ such that
   \begin{equation*}
        2 \Phi_{\kappa,N}(x) \le \Phi_{\kappa,N}(x-\delta) + \Phi_{\kappa,N}(x+\delta+2N^{-1/6}) + \varepsilon',
    \end{equation*}
    for any $N>N_{\varepsilon'}$, which implies
    \begin{equation} \label{eq:T_t_continuity_1}
        \Phi_{\kappa,N}(x) - \Phi_{\kappa,N}(x-\delta) \le \Phi_{\kappa,N}\big(x+\delta+2N^{-1/6}\big) - \Phi_{\kappa,N}(x) + \varepsilon'. 
    \end{equation}
    Consider a non-negative integer $k$ such that
    \begin{equation*}
       x+k\delta+k(k+1)N^{-1/6} \le 1 < x+(k+1)\delta+(k+1)(k+2)N^{-1/6}.
    \end{equation*}
    Then, iterating \eqref{eq:T_t_continuity_1} we obtain
    \begin{align} \nonumber
            \Phi_{\kappa,N}(x)-\Phi_{\kappa,N}(y) &\le \Phi_{\kappa,N}\left(x+k\delta+k(k+1)N^{-\frac16}\right)-\Phi_{\kappa,N}\left(x+(k-1)\delta+k(k-1)N^{-\frac16}\right) + k\varepsilon'
            \\
            &\le \Phi_{\kappa,N}(1) - \Phi_{\kappa,N}\bigg(1-2\delta-(4k+2)N^{-1/6}\bigg) + k\varepsilon'. \label{eq:T_t_continuity_2}
    \end{align}
    Next, we estimate the term $\Phi_{\kappa,N}(1) - \Phi_{\kappa,N}\big(1-2\delta-(4k+2)N^{-1/6}\big)$. By \Cref{p.energystab} and \Cref{prop:f_limit_g+T} we know that there exists $x_0=x_0(q,\kappa,a)>0$ such that 
    \begin{equation*}
    \begin{split}
        \mathcal{F}_\kappa(x_0) = \Phi_{\kappa,N}(1) + g(x_0-1) = \lim_{N \to +\infty} \inf_{y\in [0,1]} \left\{ \Phi_{\kappa,N}(y) + g(x_0-y) \right\},
    \end{split}
    \end{equation*}
    where $g(y)= \logq y^2/2 $.
    This implies that for any fixed $\varepsilon''$ we can pick $N_{\varepsilon''}$ such that, for all $N>N_{\varepsilon''}$
    \begin{equation*}
        -\varepsilon'' \le \inf_{y\in [0,1]} \left\{ \Phi_{\kappa,N}(y) + g(x_0-y) \right\} - \Phi_{\kappa,N}(1) - g(x_0-1) \le \varepsilon''.
    \end{equation*}
    Then, for any $y\in[0,1]$ we have
    \begin{equation}\label{eq:T_t_continuity_3}
        0\le \Phi_{\kappa,N}(1)- \Phi_{\kappa,N}(y) \le \varepsilon'' + g(x_0-y)-g(x_0-1)  \le \varepsilon'' + \beta_q(1-y),
    \end{equation}
    where $\beta_q:=\eta_q\cdot x_0$. Combining the estimates \eqref{eq:T_t_continuity_2}, \eqref{eq:T_t_continuity_3} we arrive at the bound
    \begin{equation} \label{eq:T_t_continuity_4}
       0\le \Phi_{\kappa,N}(x)-\Phi_{\kappa,N}(y) \le \varepsilon'' + 2\beta_q\delta+(4k+2)\beta_q\cdot N^{-1/6} + k \varepsilon',
    \end{equation}
    which holds for any $N>\max\{ N_{\varepsilon'}, N_{\varepsilon''} \}$. It is now clear that the right-hand side of \eqref{eq:T_t_continuity_4} can be made arbitrarily small, since $k<2/\delta$ and $\varepsilon', \varepsilon''$ are independent of $\delta$. Moreover, we can also allow $|x-y|<\delta$ using the fact that $\Phi_{\kappa,N}$ is non-decreasing. This completes the proof. 
\end{proof}

\end{document}
